\documentclass[11pt]{article}

\usepackage[margin=1in]{geometry}
\usepackage[T1]{fontenc}
\usepackage[utf8]{inputenc}
\usepackage{lmodern}
\usepackage{booktabs}
\usepackage{tabularx}
\usepackage{array}
\usepackage{enumitem}
\usepackage{hyperref}
\usepackage{xcolor}
\usepackage{longtable}

\hypersetup{
    colorlinks=true,
    linkcolor=blue,
    urlcolor=blue
}

\title{\textbf{NanoBubbleEval: Data Collection, Curation, and Benchmark Design}\\
\large Supervisor-Facing Project Summary}
\author{Anonymous Authors}
\date{\today}

\begin{document}
\maketitle

\section*{Project Overview}
NanoBubbleEval is a literature-derived benchmark project for evidence-grounded scientific information extraction. The project was developed in three stages. First, we built a large scholarly warehouse from public metadata and abstract sources. Second, we filtered that warehouse into a high-precision nanobubble core. Third, we curated benchmark-ready subsets for structured extraction, numerical grounding, and evidence attribution.

\section*{Data Collection Workflow}
The collection process was intentionally API-first and reproducible.
\begin{enumerate}[leftmargin=1.7em]
    \item We began with curated seed records from scientific literature and ClinicalTrials.gov.
    \item We expanded the corpus using structured APIs: PubMed / NCBI Entrez, Europe PMC, OpenAlex, CrossRef, Semantic Scholar, and ClinicalTrials.gov.
    \item We used query families spanning nanobubble-core terms, adjacent ultrasound and delivery terms, and broader nanoparticle fields needed for extraction coverage.
    \item We aggregated metadata such as title, authors, year, venue, DOI, PMID/PMCID when available, abstract, citation counts, and source provenance.
    \item We deduplicated records in a strict order: DOI, PMID/PMCID, normalized title, then URL.
    \item We assigned tier labels A, B, and C, then split Tier A into A1 and A2 for precision tightening.
    \item We formed benchmark candidate sets and gold-annotation pools from the curated subset.
\end{enumerate}

\section*{APIs and Sources}
\begin{itemize}[leftmargin=1.7em]
    \item \textbf{PubMed / NCBI Entrez}: high-volume biomedical abstracts and metadata.
    \item \textbf{Europe PMC}: abstract records, open-access signals, and citation counts.
    \item \textbf{OpenAlex}: graph expansion, references, citations, and broad scholarly metadata.
    \item \textbf{CrossRef}: DOI-centric metadata and venue discovery.
    \item \textbf{Semantic Scholar}: citation-graph expansion around high-priority works.
    \item \textbf{ClinicalTrials.gov}: translational and clinical records.
    \item \textbf{Curated seeds}: a small bootstrap set to anchor expansion.
\end{itemize}

\section*{What Each Record Contains}
Records may contain:
\begin{itemize}[leftmargin=1.7em]
    \item title, authors, year, venue, DOI, URL
    \item PMID and PMCID when available
    \item abstract or short summary
    \item source API and source platform
    \item document type and access level
    \item tier label and nanobubble-vs-nanoparticle label
    \item application category
    \item likely extraction-field flags for size, zeta potential, stability, payload, loading efficiency, and release profile
    \item relevance and extraction-priority scores
    \item notes and placeholders for evidence spans
\end{itemize}

\section*{Dataset Structure}
\begin{tabularx}{\textwidth}{>{\bfseries}p{3.5cm}X}
\toprule
Warehouse & 49,182 deduplicated records. \\
High-precision core & 7,785 records, split into A1 and A2. \\
Benchmark candidates & 500-record and 1,000-record curated subsets. \\
Gold annotation pools & 200-record and 500-record ranked candidate pools. \\
Task packages & Structured extraction, numerical grounding, and evidence attribution files. \\
\bottomrule
\end{tabularx}

\section*{Quantitative Summary}
\begin{table}[h]
\centering
\begin{tabular}{lrr}
\toprule
Metric & Count & Comment \\
\midrule
Warehouse size & 49,182 & deduplicated \\
High-precision core & 7,785 & A1 + A2 \\
A1 & 3,573 & explicit nanobubble core \\
A2 & 4,212 & adjacent ultrasound/gas-carrier \\
Benchmark candidate 500 & 500 & balanced \\
Benchmark candidate 1000 & 1,000 & balanced \\
Gold annotation v2 & 200 & manual review pool \\
Gold annotation v3 & 500 & ranked candidates \\
\bottomrule
\end{tabular}
\end{table}

\begin{table}[h]
\centering
\begin{tabular}{lrr}
\toprule
Field & Count & Note \\
\midrule
Size & 12,838 & strong coverage \\
Zeta potential & 3,350 & weakest field \\
Stability & 9,964 & strong \\
Payload & 27,369 & strongest field \\
Loading efficiency & 9,657 & improved but still limited \\
Release profile & 12,462 & strong in delivery papers \\
\bottomrule
\end{tabular}
\end{table}

\section*{Strengths}
\begin{itemize}[leftmargin=1.7em]
    \item Large and reproducible literature warehouse.
    \item Clear distinction between warehouse, core subset, and benchmark subsets.
    \item Explicit evidence-oriented fields and scores.
    \item Broad coverage across nanobubble, delivery, imaging, and nanoparticle-adjacent literature.
\end{itemize}

\section*{Current Limitations}
\begin{itemize}[leftmargin=1.7em]
    \item Tier C is intentionally broad and includes many nanoparticle-adjacent papers.
    \item ClinicalTrials.gov coverage is still small.
    \item Zeta potential remains underrepresented.
    \item Some records are metadata-rich but not full-text rich.
    \item The A1 precision proxy is heuristic and should not be treated as gold annotation.
\end{itemize}

\section*{Why This Matters for NeurIPS}
The project supports scientific NLP evaluation on:
\begin{itemize}[leftmargin=1.7em]
    \item structured extraction,
    \item numerical grounding,
    \item evidence attribution,
    \item cross-document comparison,
    \item and domain-specific literature reasoning.
\end{itemize}
That combination makes the dataset useful both as a benchmark resource and as a case study in large-scale scientific information extraction.

\section*{Questions My Supervisor May Ask}
\begin{itemize}[leftmargin=1.7em]
    \item \textbf{Why 49K records but only 7.7K in the core?} Because the warehouse is broad, while the core is deliberately filtered for nanobubble precision.
    \item \textbf{What is the difference between warehouse and benchmark?} The warehouse is the full corpus; the benchmark is a curated evaluation subset.
    \item \textbf{Why is A1 more trustworthy?} It requires explicit nanobubble lexical evidence and excludes most broad nanoparticle-only records.
    \item \textbf{What is inside one record?} Metadata, provenance, tiering, extraction flags, and benchmark-ready placeholders.
    \item \textbf{What comes next?} Manual annotation, baseline systems, and final evaluation reporting.
\end{itemize}

\section*{Next Steps}
\begin{enumerate}[leftmargin=1.7em]
    \item Finalize annotation guidelines.
    \item Manually verify the gold pools.
    \item Run baseline extractors on the benchmark files.
    \item Report extraction, numerical grounding, and evidence attribution metrics.
\end{enumerate}

\end{document}
